% !TEX root = 00_dmlr_compel.tex
\section{Discussion and Limitations}\label{sec:discussion}

\textsc{Compel} was designed to address a fundamental challenge in large-scale language model training: identifying and filtering high-quality data from internet-scale corpora without incurring the cost of large models, supervision, or inference.
While individual gains from \textsc{Compel} filtering are modest (+0.4--1.1 points), our results demonstrate that these gains are consistent across datasets and model sizes, achieved with negligible computational overhead.

Prior work shows that small, reproducible accuracy improvements can have outsize benefits at scale \citep{marion2023less, ankner2024perplexed, tirumala2023d4}.
For example, a 1-point improvement in accuracy at trillion-token scale can eliminate the need for hundreds of thousands of additional training steps or significantly reduce reliance on larger, more costly models.

By offering a simple, domain-agnostic filter that drops into existing pipelines, \textsc{Compel} provides a practical path toward higher-quality data at minimal cost---helping bridge the gap between dataset scale and dataset quality.

Our approach has several limitations.
First, compression thresholds were manually tuned based on observed distributions from reference corpora, which may not generalize optimally across diverse datasets or content domains.
While these thresholds proved broadly effective, our analysis did not systematically explore their sensitivity or optimality due to computational constraints.
Future work could automate this process by using adaptive threshold selection techniques that optimize filtering performance on small validation sets or leverage unsupervised clustering over compression distributions.

Second, applying a fixed, global compression threshold across entire corpora inherently ignores variability across different content types or domains.
More granular, adaptive filtering---especially tailored to specific domains such as code repositories, news, or social media content---could yield further performance gains.

Third, we positioned \textsc{Compel} as a final refinement step in existing filtering pipelines.
Although this demonstrated its complementary nature and consistent improvements, we believe compression-based filtering may provide greater value when applied earlier in data pipelines.
Early-stage filtering can reduce the amount of data subject to more expensive processing, making full pipelines more efficient.

Finally, the inherent simplicity of compression ratio as a scalar metric means it cannot capture semantic or nuanced linguistic content distinctions, consistent with the compressor-relative caveat of Section~\ref{sec:theory}.
Hybrid approaches that combine compression-based signals with lightweight semantic or lexical features could address this shortcoming, refining dataset quality further without incurring significant computational overhead.

Together, these limitations point to several promising directions: (1) adaptive thresholding based on unsupervised signal analysis, (2) domain-specific calibration, (3) early-stage integration into multi-pass pipelines, and (4) hybrid filters that combine compression with complementary lightweight signals.
